\section{Groups, Subgroups and Cosets}

{\bf Abstract Group}\\
A set of elements $S$ forms a group if to any two elements $x$ and $y$ of $S$ taken in a particular order there is associated a unique third element of $S$, called their product and denoted by
$xy$, which satisfies the following laws

\begin{itemize}
\item[~~(M2)] For any three elements $$x, y \mbox{ and } z: ~~ (xy)z = x(yz).$$ 
\item[~~(M3)] There is an element called the {\it neutral element } and denoted by $e$ which has the property that $xe = ex = x$ for all elements $x$.
\item[~~(M4)] Corresponding to each element $x$ there is an element $x^{-1}$ called the {\it inverse } of $x$ which has the property that 
\[ xx^{-1} = x^{-1} x = e.\]
\end{itemize}

{\bf Subgroups }\\
A non-empty subset $H$ of a group $G$ is a subgroup of $G$ if: 
\begin{itemize}
\item[(a)] Given any ordered pair $(h_1, h_2)$ of elements of $H$ the product $h_1 h_2$ as defined in $G$ is in $H$.
\item[(b)] The neutral element of $G$ is in $H$.
\item[(c)] Given any element $h$ in $H$, the inverse $h^{-1}$ as defined in $G$ is in $H$. 
\end{itemize}
This a subset of elements of $G$ is a subgroup, if and only if it forms a group, under the same group product. \\

Note that $(b)$ is not needed, since if $h\in H, h^{-1} \in H$ by $(c)$ and so $h h^{-1} = e \in H$ by $(a)$. \\

\begin{theorem}
A necessary and sufficient condition for a non-empty subset $H$ to be a\\ subgroup of a group $G$ is that $gh^{-1} \in H$ for all $g, h \in H$.
\end{theorem}

\begin{itemize}
\item[(1)] Necessary: If $H$ is a subgroup $h^{-1}\in H$ and so $gh^{-1} \in H$.
\item[(2)] Sufficient: If $g$ is any element of $H$ we have $gg^{-1} \in H$, i.e, $e\in H$. Hence $eg^{-1} = g^{-1}\in H$, i.e, $H$ contains inverses. Thus if $g,h \in H$, so is $h^{-1}$ and hence
$g(h^{-1})^{-1} \in H$, i.e. $gh \in H$.
\end{itemize}

Given a group $G$ and a subgroup $H$ it is possible to decompose the whole group into subsets `parallel' to $H$, two elements being in the same subset if their `difference' is in $H$. This concept is 
better adapted in the product notation where the difference between $g$ and $k$ is expressed as $k^{-1}g$. We therefore decompose group $G$ into subsets such that $g$ and $k$ are 
in the same subset if $k^{-1}g\in H$. These subsets are known as cosets and we have a true decomposition into mutually exclusive subsets because $k^{-1}g\in H$ sets up an equivalence
relation $gRk$ if $k^{-1}g\in H$.

\begin{theorem}{Equivalence Classes of $G$ relative to $H$}
The relation defined by $gRk$ above is an equivalence relation.
\end{theorem}

\begin{itemize}
\item[(1)] Reflexive: For any $g\in H$ we have $gRg$ because $g^{-1}g = e \in H$ since $H$ is a subgroup.
\item[(2)] Symmetric: If $gRk$ then $k^{-1}g \in H$. Hence its inverse is in H, i.e. $g^{-1}k\in H$ and so $kRq$.
\item[(3)] Transitive: If $gRk$ and $kRl$ then $k^{-1}g \in H$ and $l^{-1}k \in H$. Hence $(l^{-1}k)(k^{-1}g) = l^{-1}g \in H$ and so $gRl$.
\end{itemize}

These mutual exclusive classes are called left cosets of $G$ relative to $H$. So if $gRk$ then $kRg$ and $g^{-1}k \in H$ or $k = gh$ for some $h\in H$. Hence the coset containing g is
precisely the complex $gH$, the set $\{gh\}$ for all $h\in H$.

\begin{theorem}{Left Cosets}
The left cosets of $G$ relative to $H$ are the complexes $gH,~g\in G$. Any left coset may be expressed in this form for any $g$ in it. Any two cosets are either the same of have
no element in common.
\end{theorem}

We may similarly define \it  right cosets \rm of $G$ relative to $H$ as the equivalence classes under the relation given by $gk^{-1}\in H$. They are the complexes $Hg$.\\

We normally will use left cosets and therefore will just refer to them as cosets. 

\section{Invariant Subgroups}

\begin{definition}
$H$ is an invariant subgroup of $G$ if it is a subgroup and if the left and right cosets of $H$ in $G$ coincide; i.e., if 
\[ gH = Hg \quad \forall g \in G.\]
\end{definition}

Invariant subgroups are often called normal subgroups or self-conjugate subgroups. 

If the left coset $gH$ is also a right coset, then it contains  $g (= eg)$ and so must be the right coset $Hg$. Hence is every left coset is a right coset (or vice versa) $H$ is invariant.\\

Let us show that this condition is sufficient.
\begin{theorem}
If $H$ is invariant in $G$ then $(gH) (g'H) = g g' H$ for all $g,g' \in G$. 
\end{theorem} 

\begin{proof}
\begin{eqnarray*}
 gH g'H &=& g (Hg') H = g(g'H)H \mbox{ since } Hg' = g'H  \mbox{ as H is invariant,}\cr
 &=& gg'H^2 = gg'H \mbox{ since } H^2 = H \mbox{ as } H \mbox{ is a subgroup.}
\end{eqnarray*}
\end{proof}
Note that is would have been sufficient to prove that \[ (gH)(g'H) \subseteq gg'H,\] 

\begin{theorem}
A subgroup $H$ is an invariant subgroup of $G$ is and only if $g^{-1}Hg = H$ for all $g\in G$. 
\end{theorem}

\section{Conjugacy}
\begin{definition}
The elements $x$ and $g^{-1} x g$, where $x$ and $g$ are elements of a group $G$, are called conjugate elements in $G$ : in other words $x$ and $y$ are conjugate if and only if there exists
and element $g \in G$ such that $y = g^{-1} x g$.
\end{definition}

\begin{theorem}
The relation of conjugacy between elements of any group $G$ is an equivalence relations. 
\begin{itemize}
\item[]{\it ~~~~~Reflexive.} Since $x = e^{-1} x$, $x$ is conjugate to itself.
\item[]{\it ~~~~~Symmetric.} If $ y = g^{-1} x g$ we have $x = g y g^{-1} = (g^{-1})^{-1} y g^{-1}$. 
\item[]{\it ~~~~~Transitive.}  If $  y = g^{-1} x g$ and  $z = h^{-1} y h$ we have \[ z = h^{-1}g^{-1} x g h = (gh)^{-1} x (gh).\] 
\end{itemize}
\end{theorem}
It follows that the relation divides the elements of $G$ into mutually exclusive classes; these are called {\it conjugacy classes in G}.

\begin{theorem}
If $H$ is a subgroup of $G$ then $g^{-1} H g$ is a subgroup for any element $g\in G$. 
\end{theorem}
\begin{proof}
For \[(g^{-1} H g)^2 = (g^{-1} H g g^{-1} H g) = (g^{-1} H H g) = (g^{-1} H g) \mbox{   since }  H^2 = H\] as $H$ is a subgroup. Also
\[  (g^{-1} H g)^{-1} = g^{-1} H^{-1} g = g^{-1} H g  \mbox{   since } H^{-1} = H.\]  
\end{proof}

We say that $H$ and $g^{-1} H g$ are {\it conjugate subgroups} in $G$. 

\begin{theorem}
The relation of conjugacy between subgroups of $G$ is an equivalence relation.
\end{theorem}

\begin{theorem}
Normal subgroups of $G$ are composed of all elements of $g$ which are common to the set of conjugate subgroups of $G$. 
\end{theorem}
\begin{proof}
Suppose the set of conjugate subgroups of $G$ are
\[ H_1, H_2, \hdots H_i.\]
The elements which are common to these subgroups form a group, since with any two of them contained in $H_i$ their product is also contained in $H_i$. 
Also, the transformations of $H_1$ by elements in $G$ are the same conjugate subgroups
\[ H_1, g_2^{-1}H_1 g_2, \hdots, g_i^{-1} H_1 g_i, \] the group of common elements, which is part of any one of them as it is of $H_1$, is identical with its 
transforms and hence a normal subgroup of $G$. 
\end{proof}
A group of all elements which are common to conjugate subgroups, is called their {\bf greatest common subgroup} and a greatest common subgroup 
is normal in the group. 

\begin{theorem}
$H$ is an invariant subgroup of $G$ if the conjugacy class of subgroups of $G$ which contains $H$ consists of just $H$ itself: i.e., if the only subgroup
conjugate to $H$ is $H$. 
\end{theorem}

\section{Quotient groups}
\begin{theorem}
If $H$ is an invariant subgroup of $G$, the cosets of $H$ in $G$ form a group under the group product given by
\[ (gH) (g'H) = gg'H.\]
\end{theorem}
The product is well defined and unique, and is also a coset. \\

It is associative, since 
\[ ((gH)(g'H)) (g'' H) = gg'g''H = (gH)( (g'H)(g''H)).\]

The coset $H = (eH)$ is a unity since
\[ (gH) (eH) = geH =gH \mbox{   and  } (eH)(gH) = egH = gH.\]

The coset $gH$ has an inverse $g^{-1}H$, since
\[ (gH)(g^{-1}H) = gg^{-1}H = eH = H \mbox{   and   } (g^{-1}H)(gH) = H\] similarly.


The quotient group formed by the cosets of $H$ and $G$ is called the {\it Quotient Group of $H$ in $G$}, or sometimes the {\it Factor Group of $G$ relative to $H$}.

\section{Quotient Group Associated with a Homomorphism}

As we have already hinted, the subjects of invariant subgroups and homomorphisms of groups are closely connectedl there is in fact a 1-1 correspondence between subgroups of a given group $G$ and the essentially different 
homomorphism which have $G$ ax the object space, so that the study of either tops is similar to that of the other. \\

Suppose that $H$ is an invariant subgroup of $G$, so that the quotient group $G/H$ is defined. Consider the mapping $\theta: G \goes G/H$ defined by $g\theta = gH$, the coset containing $g$, which is of course an element of $G/H$.

\begin{theorem}
The mapping $\theta: G \goes G/H$ defined by $g\theta = gH$ is a homomorphism.
\end{theorem}

\begin{proof}
\begin{eqnarray*}
(gg')\theta &=& gg'H = gg' H H\cr
                &=& gH g'H \cr
                &=& (g\theta) (g'\theta)
\end{eqnarray*}
\end{proof}

\begin{corollary}
The kernel $\theta$ is $H$ and the image is $G/H$, so that $\theta$ is an epimorphism.
\end{corollary}

The homomorphism $\theta$ is called the Natural Homomorphism or the Canonical Homomorphism associated with $H$. It is essential that $H$ be an invariant subgroup, for otherwise $G/H$ is not defined and we cannot speak of the product of the cosets $(gH)(g'H)$.

Any homomorphism $\theta: G\goes G'$ is closely related to the natural homomorphsim associated with an invariant subgroup of $G$. 

\begin{theorem}
The kernel of $\theta: G \goes G'$ is an invariant subgroup of $G$. 
\end{theorem}

If $h$ is in the kernel $H$ and $g$ is any element of $G$, we have
\begin{eqnarray*}
(g^{-1} h g)\theta  &=& (g^{-1}\theta) (h\theta) (g\theta)\cr
&=& (g\theta)^{-1}) f (g\theta), \mbox{  where $f$ is the neutral element of $G'$}\cr
&=& (g\theta)^{-1}(g\theta) = f
\end{eqnarray*}

\begin{theorem}
The mapping defined by $gH \goes g\theta$ is a well-defined mapping of $G/H$ onto $G\theta$.
\end{theorem}
If $g_1$ and $g_2$ are in the same coset, then $g_1\theta = g_2\theta$. If $g_1$ and $g_2$ are in the same coset, then $g_2^{-1}g_1 \in H$. Hence
$(g_2\theta)^{-1}(g_1\theta) = f$, the neutral element in $G'$, since $H$ is in the kernel of $\theta$. Therefore $g_2\theta = g_1\theta$. 

Also  any element of $G\theta$ is the image of some $g\in G$ and so the mapping is onto $G\theta$, although not necessarily onto $G'$. 

\begin{theorem}
The mapping defined by $gH \goes g\theta$ is an isomorphism between $G/H$ and $G\theta$. 
\end{theorem}
We have that 
\begin{eqnarray*}
(g_1 H) (g_2 H) &=& g_1 g_2 H, \cr
(g_1\theta) (g_2\theta) &=& (g_1 g_2)\theta, 
\end{eqnarray*}
because $\theta$ is a homomorphism and $g_1 g_2 H \goes (g_1 g_2)\theta$. 

Since this mapping is onto $G\theta$, it is also an isomorphism provided it is 1-1, i.e., provided the kernel is the neutral element of $G/H$. Suppose
$gH \goes f$. Then $g\theta \goes f$ and so $g\in H$, so that $gH = H$, the neutral element of $G/H$. \\

Notice that the sets of elements of $G$ which map into the particular elements of $G\theta$ are the cosets of $G$ relative to the kernel $H$. \\

We have proved that to any invariant subgroup $H$ of $G$ there is a natural homomorphism from $G$ onto $G/H$, and conversely that any homomorphism of G is  isomorphic to
a natural homomorphism associated with its kernel. Therefore there is a 1-1 correspondence between the invariant subgroups and the essential different homomorphisms of $G$. By 
essentially different homomorphisms we mean those having different kernels: two homomorphisms having the same kernel are both equivalent to the same natural homomorphism and therefore exhibit the same structure.\\ 

The image space $G'$ is largely irrelevant : the subgroup $G\theta$ of $G'$ is isomorphic to $G/H$ but $G'$ may be larger. Also there is no reason why $G\theta$ should be invariant in $G'$. \\

{\bf Lattice Isomorphism Theorems}
 
\begin{theorem} [First Isomorphism Theorem:] Let $\phi$ be a group homomorphism from $G$ to $G'$. Then the mapping from $G/Ker(\phi)$ to $G\phi$ given by 
$g Ker(\phi) \goes g\phi$ is an isomorphism.
\end{theorem}
\begin{example}Consider the homomorphism $\phi: \Z \goes \Z_4$ such that $x\phi = x\mod4$.
In this case let $G$ be the group of integers (under addition) \[ G = \{..., -4, -3, -2, -1, 0, 1, 2, 3, 4, ...\}\] 
and \[ G' = \Z_4 = \{ 0, 1, 2, 3\}. \]
Also \[ Ker(\phi) = \{ ...-8, ,-4, 0, 4, 8, ...\} = \langle4\rangle \mbox{   (the cyclic subgroup generated by 4)} \] 
So $G\phi = \{ 0, 1, 2, 3\}$. Note that $G\phi$ doesn't have to be the entire group $G'$. In this example it is the case that $G\phi = G'$.
Then \[ G/ Ker(\phi) = \{ 0 + \langle4\rangle, 1 + \langle4\rangle, 2 + \langle4\rangle, 3+\langle4\rangle \} \mbox{   (four cosets)}\]
Notice that we can create a mapping $\overline{\phi}$ that mapps $G/Ker(\phi)$ into $G\phi$ as follows: 
\[ \overline{\phi}: G/Ker(\phi) \goes g\phi, \] where
\begin{eqnarray*}
\overline{\phi}(0+\langle4\rangle) &=& (0)\phi = 0\cr
\overline{\phi}(1+\langle4\rangle) &=& (1)\phi = 1\cr
\overline{\phi}(2+\langle4\rangle) &=& (2)\phi = 2\cr
\overline{\phi}(3+\langle4\rangle) &=& (3)\phi = 3
\end{eqnarray*}
\end{example}


\begin{theorem} [Second Isomorphism Theorem:] The Second Isomorphism Theorem: If  $S$
 is a subgroup of $G$  and $N$  is a normal subgroup of $G$ , then the product $SN$ is a subgroup, and the quotient 
$SN / N$  is isomorphic to  $S / (S \cap N$.\end{theorem}

\begin{theorem}[Third Isomorphism Theorem:] Normal Subgroups of Quotients: If $N$ and $H$ are normal subgroups of $G$ (where $N \subseteq H$), then 
$H/N$ is  is a normal subgroup of $G/N$. Furthermore, the second quotient is 
isomorphic to the primary quotient: \[ (G / N )/  (H / N)  \cong G / H. \]
\end{theorem}
The Third Isomorphism Theorem is analogous to fractions:
\[ {G / N \over H / N} \cong {G\over H}.\]

\begin{theorem} [Fourth Isomorphism Theorem:]
There is a direct, structure-preserving correspondence between the subgroups of a quotient group $G / N$
 and the subgroups of $G$
 that contain $N$.
\end{theorem}

\begin{theorem}
The mapping $\phi: \Z\goes\Z_n$ defined by $m\phi = m\mod n$ is a homomorphism with kernel $\langle n \rangle$. By the 
First Isomorphism Theorem, $\Z/\langle n\rangle \cong \Z_n$. 
\end{theorem}

If $\phi$ is a homomorphism from a finite group $G$ to $G'$, then $\|G\phi\|$ divides $\|G\}$ and $\|G'\|$. This follows from 
the First Isomorphism Theorem and the fact that $G\phi$ is a subgroup of $G'$ and Lagrange's Theorem. 

\begin{lemma}
The inner automorphism $\theta_x$ given by $g \goes x^{-1} g x$ is the identity automorphism if and only if $x\in C$ where $C$ is the center. 
\label{lemma1}\end{lemma}
If $x\in C$, $x^{-1} g x = x^{-1} x g = g$ and so $\theta_x = I$, the identity. 

If $\theta_x = I, x^{-1} g x = g \forall g \mbox{   and so  } x\in C$.

\begin{lemma}
The mapping $x \goes \theta_x$ of $G$ into the group of automorphisms of $G$ is a homomorphism
\end{lemma}
We already have that $\theta_{xy} = \theta_x\theta_y$ and so the result follows. 

By Lemma(\ref{lemma1}) the kernel of this homomorphism is the center $C$. The image is the subgroup of inner automorphisms. Therefore we obtain
\begin{theorem}
The group of inner automorphisms of a group $G \cong G/C$, where $C$ is the center of $G$.
\end{theorem}

\section{  Relationship between Normal Subgroups and Their Images}
Consider a group $G$ that has a normal subgroup $N$ and $H$ and that $H \subseteq N$. 
\begin{theorem}
Because $N\supseteq H$, both $N$ and $H$ are normal subgroups of $G$. Thus $G/N$ is the homomorphic image of $G/H$ under the natural homomorphsim
$\pi: G/H \goes G/N$ defined by 
$(xG)\pi = xN$. 
\end{theorem}
\begin{proof}

1) Define a function $\phi: G\goes G/N$ by
\[ x\phi = xN.\]
This is the standard canonical projection map. It is a surjective group homomorphism with $ker(\phi) = N$. 

2) We want to construct a map $\overline{\phi}: G/H \goes G/N$.
Define $\overline{\phi}$ by:
\[ (xH) \overline{\phi} = x\phi = xN.\]

3) We must show that if $xH = yH$, then $(xH)\overline{\phi} = (yH)\overline{\phi}.$
\begin{eqnarray*}
xH &=& yH \goes y^{-1}x \in H.\cr
H &\subseteq& N \mbox{  given}\cr
y^{-1}x &\in& N\cr
xN &=& N\cr
(xH)\overline{\phi} = (yH)\overline{\phi}\cr
\end{eqnarray*}
The map is well-defined because $ H \subseteq ker(\phi)$. 
\end{proof}

4) $\overline{\phi}$ is a Homomorphism:
\[ (xH) (yH) \overline{\phi} = (xyH)\overline{\phi} = xyN = (xN)(yN) = (xH)\overline{\phi} (yH)\overline{\phi}.\]

5) $\overline{\phi}$ is Surjective:
Let $yN \in G/N$. 
Since $y\in G$, we can form the coset $yH \in G/H$:
\[ (yH)\overline{\phi} = yN.\] Every element in the co-domain has a pre-image
 
6) Apply the First Isomorphism Theorem
By the First Isomorphism Theorem, the domain modulo the kernel is isomorphic to the image

\[ (G/H)/ker(\overline{\phi}) \cong Im(\overline{\phi}) = G/N .\]
Since $G/N$ is isomorphic to a quotient group of $G/H$, it is by definition a homomorphic image of $G/H$.\qed \\